\subsection{Lines and closed hyperplanes}

PROPOSITION 2. --- Every Hausdorff topological vector space E of dimension 1 over a non-discrete valued division ring K is isomorphic to \(K_{s}\) ; in fact, for every \(a \neq 0\) in E, the mapping \(\xi \mapsto \xi a\) of \(K_{s}\) on E is an isomorphism (in other words every linear mapping of \(K_{s}\) on E is an isomorphism).

As the mapping \(\xi \mapsto \xi a\) of \(K_{s}\) on E is bijective and continuous (I, p.~1, def. 1), it is sufficient to show that it is bicontinuous. Let \(\alpha\) be a real number \(>0\) , we show that there exists a neighbourhood V of 0 in E such that if \(\xi a \in V\) then \(|\xi| < \alpha\) . As K is not discrete, there exists an element \(\xi_0 \in \mathbf{K}\) such that \(0 < |\xi_0| < \alpha\) ; but, as E is Hausdorff, there is a neighbourhood V of 0 such that \(\xi_0a\) does not belong to V. We can suppose that V is balanced (I, p.~7, prop. 4). But then if \(\xi a \in \mathbf{V}\) and \(|\xi| \geqslant |\xi_0|\) we have \(|\xi_0\xi^{-1}| \leqslant 1\) , and \(\xi_0a = (\xi_0\xi^{-1})(\xi a) \in \mathbf{V}\) ; since this last statement is false we see that \(\xi a \in \mathbf{V}\) implies \(|\xi| < |\xi_0| < \alpha\) . This completes the proof.

COROLLARY 1. --- In a Hausdorff topological vector space E, over a non-discrete valued division ring K, every vector subspace D of dimension 1 is isomorphic to \(K_{s}\) .

COROLLARY 2. --- Let E be a topological vector space over a non-discrete valued division ring. Every vector subspace D (of dimension 1) which is the algebraic complement of a closed homogeneous hyperplane H is also the topological complement of H.

In D, the set \(\{0\}\) is closed, being the intersection of D and the closed set H; D is therefore Hausdorff. But as E/H is also Hausdorff, the canonical mapping of D on E/H, which is linear, is also an isomorphism by prop. 2, from which the conclusion follows (GT, III, § 6.2).

THEOREM 1. --- Let E be a topological vector space over a non-discrete valued division ring. Let H be a hyperplane in E defined by the equation \(f(x) = \alpha\) where f is a linear form not identically zero. Then H is closed in E if and only if f is continuous.

The condition is evidently sufficient (GT, I, § 2.1, th. 1); we show that it is necessary. We can suppose that H is a closed homogeneous hyperplane with the equation \(f(x) = 0\) . The quotient space E/H is then a Hausdorff topological vector space of dimension 1 on K. We can write \(f = g \circ \phi\) , where \(\phi\) is the canonical mapping of E on E/H and g is a linear mapping of E/H on \(K_{s}\) ; from prop. 2, g is continuous, thus the same is true of f.

COROLLARY. --- Every continuous linear form on E that is not identically zero is a strict morphism of E on \(K_{s}\) .

Remark. --- There are examples of normed topological vector spaces over a complete non-discrete valued division ring, in which every continuous linear form is identically zero (I, p.~25, exerc. 4); in such a space therefore, every hyperplane is everywhere dense (I, p.~11, corollary).

